% Part: lambda-calculus
% Chapter: lambda-definability
% Section: pairs

\documentclass[../../../include/open-logic-section]{subfiles}

\begin{document}

\olfileid{lam}{ldf}{pai}
\olsection{जोड्या आणि पूर्ववर्ती फलन}

\begin{defn}
$M$ आणि $N$ यांची जोडी ($\tuple{M,N}$ असे लिहितात)
पुढीलप्रमाणे परिभाषित केली जाते:
\[
\tuple{M,N} \ident \lambd[f][fMN].
\]
\end{defn}

सहज अर्थाने ही जोडी म्हणजे एक फलन: ते दुसरे फलन
स्वीकारते आणि ते जोडीतल्या दोन्ही घटकांना लागू करते.
या कल्पनेवरून दोन पदे घेऊन त्यांची जोडी परत करणारे
पुढील रचनाफलन मिळते:
\[
  \fn{Pair} \ident \lambd[mn][\lambd[f][fmn]]
\]
जोडी दिल्यावर तिचे घटक पुन्हा मिळवायचे असतात.
त्यासाठी जोडी युक्तिवाद म्हणून स्वीकारून तिचा पहिला
किंवा दुसरा घटक परत करणारी दोन घटक-निवड फलने हवीत:
\begin{align*}
  \fn{Fst} & \ident \lambd[p][p(\lambd[mn][m])]\\
  \fn{Snd} & \ident \lambd[p][p(\lambd[mn][n])]
\end{align*}

\begin{prob}
  घटक-निवड फलने $\fn{Fst}$ आणि $\fn{Snd}$ का कार्य
  करतात ते स्पष्ट करा.
\end{prob}

आता जोड्यांच्या साहाय्याने पूर्ववर्ती फलन
!!{lambda define} करता येते:
\[
  \fn{Pred} \ident \lambd[n][\fn{Fst}(n (\lambd[p][\tuple{\fn{Snd}\, {p}, \fn{Succ}(\fn{Snd}\, {p})}]) \tuple{\num 0, \num 0})]
\]
$\num n\, f x$ चे $f^{n}(x)$ मध्ये न्यूनीकरण
होते हे लक्षात ठेवा. येथे $f$ हे जोडी $p$ स्वीकारून
$p$ चा दुसरा घटक आणि त्या दुसऱ्या घटकाचा
उत्तरवर्ती यांची नवी जोडी परत करणारे फलन आहे;
$x$ ही $\tuple{0,0}$ जोडी आहे. म्हणून $n=0$
असेल, तर परिणामी जोडी $\tuple{0,0}$ असते;
अन्यथा ती $\tuple{\num{n-1}, \num n}$ असते.
मग $\fn{Pred}$ या परिणामी जोडीचा पहिला घटक
परत करते.

नैसर्गिक संख्यांवरील शून्यावर थांबणारी वजाबाकी
म्हणजे $a$ ला $\fn{Pred}$ हे $b$ वेळा लागू करणे:
\[
\fn{Sub} \ident \lambd[ab][b \fn{Pred}\, a].
\]

\end{document}
